\documentclass[11pt,a4paper]{article}
\usepackage{fontspec,amsmath,amssymb,amsthm,microtype}
\setmainfont{lmroman10-regular.otf}[BoldFont=lmroman10-bold.otf,ItalicFont=lmroman10-italic.otf,BoldItalicFont=lmroman10-bolditalic.otf]
\usepackage[margin=25mm]{geometry}
\usepackage[unicode,hidelinks]{hyperref}
\hypersetup{pdfauthor={LangevinLean Project}}
\setlength{\emergencystretch}{3em}
\setlength{\parskip}{0.4em}
\setlength{\parindent}{0pt}
\allowdisplaybreaks[2]
\raggedbottom
\title{\texorpdfstring{The full Gaussian ULA entropy theorem}{The full Gaussian ULA entropy theorem}}
\author{LangevinLean Project}
\date{}
\begin{document}\maketitle
Draft mathematical exposition; these notes may contain mistakes.


For a differentiable potential with an \(L\)-Lipschitz gradient, its normalized Gibbs target satisfying LSI with constant \(\alpha>0\), and a finite-KL initial probability law, this module proves the actual ULA bound

\[\mathrm{KL}(\rho_k\mid\nu)\le e^{-\alpha hk}\mathrm{KL}(\rho_0\mid\nu)+8hdL^2/\alpha,\quad0<h\le\alpha/(4L^2).\]

The final conclusion includes finiteness of the actual divergence. All score moments, density calculus, weak equations, entropy endpoints, and numerical compatibility estimates are supplied by the preceding modules. The theorem applies also in dimension zero.



The actual-law proof begins with finite initial KL. Score exponential integrability supplies the initial gradient moment; Gaussian smoothing supplies positive \(C^1\) densities and finite Fisher energy at positive time; LSI then supplies finite KL. The bounded-test entropy identity, its initial majorant, and compatibility \(\alpha\le L\) prove the scalar one-step estimate. Induction preserves finite KL along the actual ULA iterates, so the final geometric recurrence represents their true entropies.


\section{\texorpdfstring{Handle the singleton state space and a zero smoothness parameter}{Handle the singleton state space and a zero smoothness parameter}}


On a singleton state space every probability measure is equal: any measurable set either contains the unique point and is the whole space, or is empty. Measure extensionality and probability normalization therefore identify any two laws. In particular their KL divergence is zero.



For the substantive entropy calculation one works with \(L>0\), so that the step restriction and exponential score moment have their usual positive denominators. In positive dimension, \(\alpha\le L\) and \(\alpha>0\) guarantee this positivity. The singleton case has already been resolved directly with zero KL, so it needs no division by a smoothness constant.



When the state space is a singleton, every probability law is the same point mass and KL is zero. In the nondegenerate case, the compatibility argument rules out zero smoothness when the LSI constant is positive, permitting the subsequent divisions by \(L\).


\section{\texorpdfstring{Prove finiteness at every actual ULA iterate}{Prove finiteness at every actual ULA iterate}}


Start from the supplied finite initial KL and induct on \(k\). If \(\rho_k\) has finite KL, the direct entropy/\allowbreak{}score estimate gives \(\nabla U\in L^2(\rho_k)\). For the actual positive-time Gaussian interpolation, the Lebesgue score energy is at most \(d/(2h)\), and the propagated gradient moment is finite. Thus relative Fisher information is finite. LSI with \(\alpha>0\) then gives finite KL of the output.



The endpoint law is exactly \(\rho_{k+1}\), so this argument closes the finite-KL induction. The finiteness conclusion needs only positivity of the step; its smaller quantitative restriction is used later for contraction and forcing absorption.



The induction begins with the assumed finite initial KL. Finite current KL supplies its score moment; the positive-time Gaussian transition has finite Fisher energy; LSI then gives finite next KL. Thus each actual entropy term is finite before the recurrence is iterated.



Earlier results used here: \cite[Section~3.26.3]{LangevinLean2026ULASC}.

\section{\texorpdfstring{Derive all scalar step restrictions}{Derive all scalar step restrictions}}


In positive dimension, the bounded sinusoidal tests in compatibility give \(\alpha\le L\). Since \(\alpha>0\), also \(L>0\). Multiplying the original restriction by \(4L^2>0\) gives \(4hL^2\le\alpha\le L\), hence \(hL\le1/4\). Multiplying \(\alpha\le L\) by \(h>0\) also gives \(\alpha h\le1/4\).



The final theorem splits off the singleton case before invoking this positive-dimensional argument. No separate compatibility assumption is required.



From \(h\le\alpha/(4L^2)\) and \(\alpha\le L\), derive \(hL\le1/4\) and \(\alpha h\le1/4\). These numerical consequences are checked before invoking the scalar absorption and geometric-denominator lemmas.



Earlier results used here: \cite[Section~3.26.3]{LangevinLean2026ULASC}.

\section{\texorpdfstring{Apply LSI to the actual positive-\allowbreak{}time density}{Apply LSI to the actual positive-time density}}


Finite initial KL supplies the initial gradient moment. The actual Gaussian mixture is positive C\ensuremath{{}^{1}} relative to Gibbs, has finite KL, and has finite Fisher information at every \(t>0\). Apply the original extended-valued LSI to this density. Since both sides are finite, order is preserved on passing to their real values:

\[2\alpha\,\mathrm{KL}(\rho_t\mid\nu)\le I_t.\]

The integral form identifies \(I_t\) with \(\int\|\mathcal S(t,y)\|^2d\rho_t(y)\), using the proved equality between the measurable score field and \(\nabla\log r_t\), and the equality \(r_t\nu=\rho_t\). The second version keeps the extended-Fisher notation for the entropy-production assembly.



At positive elapsed time, the constructed relative density is positive, \(C^1\), and normalized, with finite Fisher and KL. It therefore belongs to the original LSI test class, giving \(2\alpha H(t)\le I_t\) for the actual law.



Earlier results used here: \cite[Section~3.26.3]{LangevinLean2026ULASC}.

\section{\texorpdfstring{Supply the actual analytic inputs to scalar comparison}{Supply the actual analytic inputs to scalar comparison}}


Fix any finite-KL probability law \(\mu\). The score-exponential argument gives the integrability of \(\|\nabla U\|^2\) and

\[G_0\le4L^2H_\nu(\mu)/\alpha+2dL.\]

The bounded-log entropy identity proves local absolute continuity of the actual entropy \(H(t)\) on every \([s,h]\) with \(s>0\). The actual Gaussian weak equation and Young estimate give \(H^{\prime}\le-3I/4+t^2L^2G_0+2tdL^2\) almost everywhere. The preceding LSI estimate supplies \(2\alpha H\le I\).



The explicit initial majorant \(M\) is continuous at zero, satisfies \(M(0)=H_\nu(\mu)\), and bounds \(H(t)\) for sufficiently small positive \(t\). These are exactly the premises of the integrating-factor comparison for a locally absolutely continuous entropy curve, giving

\[H_\nu(\mu P_h)\le e^{-\alpha h}H_\nu(\mu)+6h^2dL^2.\]

The singleton case instead uses equality of probability laws and nonnegativity of the right side. No density path, weak equation, entropy derivative, or moment estimate is a theorem argument.



The actual entropy curve has local absolute continuity, the almost-everywhere production bound, and an upper initial majorant continuous at zero. Together with the entropy-derived score moment, these supply every premise of the scalar one-step comparison.



Earlier results used here: \cite[Section~3.26.3]{LangevinLean2026ULASC}.

\section{\texorpdfstring{Derive the actual discrete recurrence}{Derive the actual discrete recurrence}}


The law \(\rho_k\) is a probability measure by the proved Gaussian Markov construction. Its KL is finite by the induction above. Apply the single-transition estimate with \(\mu=\rho_k\). The interpolation endpoint is precisely \(\rho_{k+1}\), yielding

\[H_{k+1}\le e^{-\alpha h}H_k+6h^2dL^2.\]

entropy of the actual ula law is the real value of the actual KL of iteration of the gaussian transition. The recurrence is derived here; it is not supplied as a premise.



Use the exact interpolation endpoints to replace the scalar initial and final entropies by \(H_k\) and \(H_{k+1}\). The comparison then gives \(H_{k+1}\le e^{-\alpha h}H_k+6h^2dL^2\) without any separately assumed recurrence.



Earlier results used here: \cite[Section~3.26.3]{LangevinLean2026ULASC}.

\section{\texorpdfstring{Sum the recurrence and retain the finite-\allowbreak{}KL guard}{Sum the recurrence and retain the finite-KL guard}}


Return the already-proved output finiteness as the first component of the conclusion. In positive dimension, the step estimates give \(\alpha h\le1/4\). The scalar geometric-sum theorem applied to the actual recurrence gives

\[H_k\le e^{-\alpha hk}H_0+\frac{6h^2dL^2}{1-e^{-\alpha h}}\le e^{-\alpha hk}H_0+\frac{8hdL^2}{\alpha}.\]

The final inequality uses \(1-e^{-x}\ge3x/4\) for \(0\le x\le1/4\). In dimension zero, equality of probability measures makes the left side zero and the result follows by nonnegativity.



The final hypotheses are differentiability, globally Lipschitz gradient, the Gibbs LSI with positive constant, a positive admissible step, and finite initial KL. Probability normalization of the target is included in LSI. Every density, moment, weak-equation, and entropy-regularity property used in the recurrence has been derived above from these assumptions.



Apply the geometric recurrence estimate to the now finite sequence of true KL values. The positive denominator bound converts the forcing sum into \(8hdL^2/\alpha\), yielding the final statement under only the displayed potential, LSI, step, and initial-entropy assumptions.



Earlier results used here: \cite[Section~3.26.3]{LangevinLean2026ULASC}.

\begin{thebibliography}{ULA-LSI}
\bibitem[ULA-SC]{LangevinLean2026ULASC}
LangevinLean Project. \emph{Unadjusted Langevin under Strong Convexity}. LL-01, 2026. Draft report, content revision 8d6335dce333. \url{https://odakyildiz.github.io/LangevinLean/docs/pdf/UnadjustedLangevin/FullProof.pdf}.
\end{thebibliography}

\end{document}
